투 포인터(Two pointer)는 한 배열의 어떤 부분 구간에 대한 문제를 풀 때, 두 개의 포인터를 사용하여 전체 구간을 탐색하지 않고 빠르게 문제를 해결하는 방식이다. 보통 제일 앞에 두 개의 포인터 \texttt{p}와 \texttt{q}를 잡고 두 포인터를 오른쪽으로 이동시키며 문제를 해결한다. 이 장에서는 투 포인터로 해결할 수 있는 문제가 무엇이고 어떻게 해결하는지 다룬다.

\section{가능한 구간을 찾는 문제}
다음과 같은 문제를 생각하자.

\begin{example} \label{ex: two-pointer}
    (15565번) $0$ 또는 $1$로만 이루어진 길이 $N$의 이진 배열에서 $1$의 개수가 $k$개 이상인 구간들 중 가장 짧은 구간의 길이를 출력하라.
\end{example}

만약 가능한 $O(N^2)$개의 모든 구간을 고려한다면 부분 합을 사용할 때 각 구간이 조건을 만족함을 $O(1)$에 확인할 수 있으므로 시간복잡도가 $O(N^2)$이 된다. 하지만 투 포인터 알고리즘을 사용하면 $O(N)$으로 줄일 수 있다.

이를 위해 해야 하는 중요한 관찰 두 가지가 있다.
\begin{obs}
    다음이 성립한다.
    \begin{itemize}
        \item 구간 $[s, e]$가 조건을 만족한다면 $[s, e]$를 포함하는 구간은 확인 필요가 없다. $[s, e]$보다 길이가 길기 때문이다.
        \item 구간 $[s, e]$가 조건을 만족하지 않는다면 $[s, e]$에 포함되는 구간은 확인할 필요가 없다. 자명히 $1$의 개수가 $k$개보다 작기 때문이다.
    \end{itemize}
\end{obs}
이제 고정된 $s$에 대해 살펴봐야 할 $e$의 범위를 생각해보자. 기본적으로 $e$는 증가한다고 가정한다. 만약 $[s, e]$가 조건을 만족한다면 $i>e$인 $[s, i]$는 살펴볼 필요가 없다. 따라서 $s$를 $s+1$로 증가시키자. 반대로 $[s, e]$가 조건을 만족하지 않는다면 $i<e$인 $[s, i]$는 살펴볼 필요가 없다. 따라서 $e$를 $e+1$로 증가시키자. 이렇게 되면 각 과정에서 $s$가 $1$ 증가하거나 $e$가 $1$ 증가하므로 최대 $2N$번 반복되며, $O(N)$의 시간복잡도에 문제를 해결할 수 있게 된다. 구현은 \Cref{code: two-pointer}을 참고하라.

\begin{code} \label{code: two-pointer}
\Cref{ex: two-pointer}의 정답 코드
\begin{lstlisting}
#include <bits/stdc++.h>
using namespace std;

int n, arr[101010], k;
int sum[101010], ans;

int main() {
    cin >> n;
    for (int i = 1; i <= n; i++) cin >> arr[i];
    cin >> k;

    for (int i = 1; i <= n; i++) sum[i] = sum[i - 1] + arr[i];

    int s = 1, e = 1;
    ans = n;
    while (true) {
        if (e > n) break;
        if (sum[e] - sum[s - 1] >= k) {
            ans = min(ans, e - s + 1);
            s++;
        }
        else e++;
    }

    cout << ans << "\n";

    return 0;
}
\end{lstlisting}
\end{code}

투 포인터라는 이름은 $s$와 $e$를 인덱스가 아니라 배열의 $s$번째 값과 $e$번째 값을 가리키는 포인터로 생각하여 나온 이름이다. $s$와 $e$를 포인터로 생각할 수 있는 이유는 $s$와 $e$는 증가만 하고, 이러한 특징이 포인터가 증가해 다음 값을 가리키는 것과 비슷하기 때문이다.

\section{투 포인터가 사용될 조건}
그렇다면 투 포인터가 사용될 수 있는 조건은 무엇일까? 다시 위의 관찰로 돌아가보면 다음이 성립해야 함을 알 수 있다.
\begin{obs}
    구간 $[s, e]$를 평가하는 함수 $f(x)$와(값이 클수록 좋다고 가정) 구간 $[s, e]$가 조건을 만족하는지 여부를 리턴하는 함수 $g(x)$에 대해
    \begin{itemize}
        \item 어떤 구간 $[s, e]$의 $f(x)$보다 $[s, e]$를 포함하는 구간의 $f$ 값이 더 작고,
        \item 어떤 구간 $[s, e]$의 $g(x)$가 거짓이면 $[s, e]$에 포함되는 구간의 $g$ 값이 거짓이어야 한다.
    \end{itemize}
\end{obs}

\Cref{ex: two-pointer}은 $f(x) = -(e - s + 1)$(음의 부호는 작을수록 좋기 때문에 붙었음), $g(x)=(1\textrm{의 개수} \geq k)$인 상황으로 볼 수 있다.

마지막으로 임의의 투 포인터 문제의 시간복잡도를 논의해보자. 매 과정에서 $g([s, e])$가 성립하는지 확인하고, 전체 과정은 $2N$번 이하이므로 $g([s, e])$를 계산하는 시간복잡도가 $T(n)$이면 전체 시간복잡도는 $N \cdot T(N)$이다.

\section{연습문제 A}
\begin{problem}
    \Cref{ex: two-pointer}은 부분 합을 사용하지 않아도 문제를 해결할 수 있다. 부분 합을 사용하지 않고 \Cref{ex: two-pointer}을 해결하는 코드를 작성하라.
\end{problem}

\section{연습문제 B}

\begin{problem}
    (1806번) 길이가 $N$인 수열의 연속한 부분수열들 중 수열의 원소들의 합이 $S$ 이상인 가장 짧은 부분수열의 길이를 출력하라. $N \leq 10^5$, $S \leq 10^9$이고 각 수열의 원소는 $10000$ 이하의 자연수이다.
\end{problem}

\begin{problem}
    (2470번, KOI10) $N$개의 수가 주어질 때 두 개를 골라 합을 $0$에 가장 가깝게 만드려고 한다. 그때 고른 두 수를 \textbf{아무거나} 하나 출력한다. $N \leq 10^5$이고 각 수는 절댓값이 $10^9$ 이하인 정수이며 모두 다르다.
\end{problem}

\begin{problem}
    (2473번, KOI10) $N$개의 수가 주어질 때 세 개를 골라 합을 $0$에 가장 가깝게 만드려고 한다. 그때 고른 세 수를 \textbf{아무거나} 하나 출력한다. $N \leq 5000$이고 각 수는 절댓값이 $10^9$ 이하인 정수이며 모두 다르다.
\end{problem}

\begin{problem}
    (7453번) $N$개의 수가 주어질 때 네 개를 골라 합이 $0$이 되는 방법의 수를 출력하라. $N \leq 4000$이고 각 수는 절댓값이 $2^28$ 이하인 정수이다.
\end{problem}

\begin{problem}
    (13364번*) $N$개의 점을 순서대로 연결하여 그린 $N-1$개의 선분의 집합과 $M$개의 점을 순서대로 연결하여 그린 $M-1$개의 선분의 집합을 생각하자. 각각을 경로로 하여 두 물체가 이동한다. 즉, 첫 번째 물체는 첫 번째 점에서 $N$번째 점의 방향으로, 두 번째 물체는 첫 번쨰 점에서 $M$번째 점의 방향으로 이동한다. 두 물체가 같은 속도로 이동할 때 두 물체 사이의 거리가 가장 가까울 때 그 값을 출력하라. 단, 한 물체가 먼저 마지막 점에 도달하면 물체는 소멸한다고 가정한다. $N, M \leq 10^5$이고 각 점의 $x$좌표와 $y$좌표는 $10000$ 이하의 음 아닌 정수이다.
\end{problem}

\begin{problem}
    (16598번) 연속된 출석 일수가 길수록 좋은 보상을 주는 앱이 있다. 그런데 버그가 나서 추가로 $P$일이 출석으로 인정되는 문제가 생겼다! 이 앱에 출석한 $N$일의 날짜 $d_i$가 주어질 때, 가능한 연속 출석 일수의 최댓값을 출력하라. 즉, 적절히 출석으로 인정될 $P$일을 고를 때 가능한 최대 연속 출석 일수를 구하는 것이다. $N \leq 2 \cdot 10^5$, $P \leq 2 \cdot 10^5$, $0 \leq d_1 < d_2 < \cdots < d_n \leq 10^6$이다.
\end{problem}

\begin{problem}
    (20085번*, IOI16) 자연수로 이루어진 길이 $N$의 수열 $a_i$와 범위 $[l, r]$이 주어진다. 이때 길이 $N$의 수열에서 적절히 몇 개의 수를 골라 합이 $[l, r]$의 범위에 속하게 할 수 있는지 판단하고, 가능하다면 선택되는 수열의 수의 인덱스를 오름차순으로 출력하라. 단, $r - l \leq \max a_i - \min a_i$가 보장된다.
    \begin{itemize}
        \item (서브테스크 1) $n \leq 100$, $a_i \leq 100$, $1 \leq l \leq r \leq 1000$, $a_i$는 모든 $i$에서 같다.
        \item (서브테스크 2) $n \leq 100$, $a_i \leq 1000$, $1 \leq l \leq r \leq 1000$, $\max a_i - \min a_i \leq 1$
        \item (서브테스크 3) $n \leq 100$, $a_i \leq 1000$, $1 \leq l \leq r \leq 1000$
        \item (서브테스크 4) $n \leq 10000$, $a_i \leq 10000$, $1 \leq l \leq r \leq 10000$
        \item (서브테스크 5) $n \leq 10000$, $a_i \leq 5 \times 10^5$, $1 \leq l \leq r \leq 5 \times 10^5$
        \item (서브테스크 6) $n \leq 100$, $a_i < 2^31$, $1 \leq l \leq r < 2^31$
    \end{itemize}
\end{problem}

\begin{problem}
    (20133번*) $N$개의 점이 주어질 때 서로 다른 \textit{모래시계 모양}의 개수를 출력하라. $N \leq 1500$이고 점의 $x$좌표와 $y$좌표는 절댓값이 $10^9$ 이하인 정수이며 어떤 세 점도 한 직선을 이루지 않는다.
    \begin{itemize}
        \item \textit{모래시계 모양}은 두 삼각형이 한 점을 공유하고 두 삼각형이 겹치지 않는 모양이다.
        \item 서로 다른 \textit{모래시계 모양}이라는 것은 구성하는 점과 선분이 중 하나라도 다른 것이다.
    \end{itemize}
\end{problem}

\begin{problem}
    (20922번) 길이 $N$의 수열이 주어질 때, 구간에 포함된 임의의 정수가 구간에 $k$개 이하로 포함되는 가장 긴 연속 부분 수열의 길이를 출력하라. $N \leq 2 \cdot 10^5$, $1 \leq k \leq 100$이고 수열의 각 정수는 $100000$ 이하의 자연수이다.
\end{problem}

\begin{problem}
    (20985번*) $1$부터 $N$까지 번호가 매겨진 $N$개의 눈덩이가 있다. $i$번 눈덩이는 $x_i$에 있고 처음 무게는 0이다. 이 눈덩이는 바람에 의해 평원을 구르며 크기가 커지게 된다. 정확히는 다음과 같은 방식이다.
    \begin{itemize}
        \item $Q$일 동안 정수 $w$로 표현되는 바람이 분다. 이 바람이 불면 하루가 끝나고 나서 눈덩이가 $x_i$에서 $x_i + w$로 이동한다. 또한 모든 눈덩이는 동시에 이동한다.
        \item 구간 $[a, a+1]$에 눈이 있다면 그 구간을 지나간 눈덩이의 무게가 1 증가한다. 그리고 눈이 사라진다. 만약 눈이 없다면 무게의 변화는 없다.
    \end{itemize}
    $Q$일이 지난 이후 $N$개의 눈덩이의 무게를 각각 출력하라.
\end{problem}

\begin{problem}
    (25405번*, KOI22) 레벨이 $L_i$인 $N$명의 게임 캐릭터에 대해 다음과 같은 규칙의 트레이닝을 $M$번 진행하였을 때 모든 트레이닝이 끝난 이후 각 캐릭터들의 레벨을 오름차순으로 정렬하여 출력하라.
    \begin{itemize}
        \item 한 번의 트레이닝은 레벨이 낮은 순서대로 $K$개의 캐릭터를 고르고 선택된 캐릭터들의 레벨을 $1$ 올리는 것이다. 레벨이 같은 캐릭터가 여러 명인 경우에는 아무거나 고른다.
    \end{itemize}
    \begin{itemize}
        \item (서브테스크 1) $N \leq 1000$, $M \leq 1000$, $K \leq N$, $L_i \leq 10^9$
        \item (서브테스크 2) $N \leq 10^5$, $M \leq 10^9$, $K = 1$, $L_i \leq 10^9$
        \item (서브테스크 3) $N \leq 10^5$, $M \leq 10^5$, $K \leq N$, $L_i \leq 10^9$
        \item (서브테스크 4) $N \leq 10^5$, $M \leq 10^9$, $K \leq N$, $L_i \leq 10^9$
    \end{itemize}
\end{problem}

\begin{problem}
    (26829번*) $n$개의 참-거짓 질문에 대한 답변이 모두 다른 $m$명의 사람이 있다. 그들의 답변은 $[0, 2^n)$ 사이의 숫자로 표현되며, $i$번째 비트가 $1$이라는 것은 $i+1$번째 질문에 대해 참으로 답변했음을 의미한다. $m$명을 인덱스 기준으로 $[1, p]$와 $[p+1, m]$의 두 그룹으로 나눌 때 다음 조건을 만족하도록 나누는 경우의 수를 구하여라. $n \leq 30$, $m \leq \min(2^n, 2 \times 10^6)$이다.
    \begin{itemize}
        \item 모든 $1 \leq i \leq n$에 대해 $i$번 질문에 참이라 대답한 사람이 \textbf{두 그룹 모두 존재하거나 두 그룹 모두 존재하지 않아야 한다}.
        \item 모든 $1 \leq i \leq n$에 대해 $i$번 질문에 거짓이라 대답한 사람이 \textbf{두 그룹 모두 존재하거나 두 그룹 모두 존재하지 않아야 한다}.
    \end{itemize}
\end{problem}

\begin{problem}
    (27916번*) $N$개의 아파트가 수직선 상에 있고, 각 아파트의 위치는 $A_i$로 표현된다. 각 아파트는 정확히 하나의 아파트 단지에 속해야 하고, 아파트 단지는 최소 $M$개의 아파트를 포함하는 연속된 번호를 가진 아파트의 집합이어야 한다. 다른 말로 모든 아파트 단지는 $[L, R]$로 표현될 수 있고, 이렇게 표현된다면 아파트 단지의 크기는 $A_R - A_L$이다. 다음 쿼리 $Q$개를 해결하는 프로그램을 작성하라.
    \begin{itemize}
        \item 아파트 단지의 크기를 $X_i$ 이하로 할 수 있는 아파트 단지의 구성 방법이 존재한다면 \texttt{1}을, 존재하지 않는다면 \texttt{0}을 출력하라.
    \end{itemize}
    $1 \leq A_1 < A_2 < \cdots < A_N \leq 10^9$, $0 \leq X_i \leq 10^9$이고 서브테스크마다 추가 제약 조건이 있다.
    \begin{itemize}
        \item (서브테스크 1) $N \leq 3 \times 10^5$, $M = N$, $Q = 1$
        \item (서브테스크 2) $N \leq 1000$, $M \leq N$, $Q \leq 3 \times 10^5$
        \item (서브테스크 3) $N \leq 3 \times 10^5$, $M \leq N$, $Q = 1$
        \item (서브테스크 4) $N \leq 1000$, $M \leq N$, $Q \leq 3 \times 10^5$
        \item (서브테스크 5) $N \leq 3 \times 10^5$, $M \leq N$, $Q \leq 3 \times 10^5$
    \end{itemize}
\end{problem}

\begin{problem}
    (28081번) 영역 $[0, W] \times [0, H]$를 $y$축과 평행한 $M$개의 선 $x=x_i$과 $x$축과 평행한 $N$개의 선 $y=y_i$로 나누면 각각의 직사각형 영역이 만들어진다. 이 직사각형 영역들 중 넓이가 $K$ 이하인 영역의 개수를 출력하라. $2 \leq H, W \leq 10^9$, $N, M \leq 10^5$, $1 \leq x_1 < x_2 < \cdots < x_M < W$, $1 \leq y_1 < y_2 < \cdots < y_N < H$이다.
\end{problem}

\section{연습문제 C}
\begin{problem}
    (8191번) 길이 $N$의 수열 $A_i$가 주어질 때, $M$번 $k$가 주어질 때마다 다음 쿼리를 해결하는 프로그램을 작성하라. $N \leq 10^6$, $M \leq 50$, $A_i \leq 10^9$, $k \leq 10^9$이다.
    \begin{itemize}
        \item $A_i > k$인 한 $i$를 골라 $A_i$를 $1$ 줄이고 $i > 1$일 때 $A_{i-1}$을 $1$ 늘리거나 $i < N$일 때 $A_{i+1}$을 $1$ 늘리는 연산을 사용할 수 있다. 이 연산을 통해 $A_i \geq k$인 연속된 구간의 길이를 최대화시키고자 할 때 가능한 최댓값을 출력하라.
    \end{itemize}
\end{problem}

\begin{problem}
    (10010번) $N$개의 도시가 있고 일부 도시 쌍이 양방향 도로로 연결되어 있아. 구체적으로 $1 \leq i < N$에 대해 $0$번 도시와 $i$번 도시가 연결되어 있고, $1 \leq i < N-1$에 대해 $i$번 도시와 $i+1$번 도시가 연결되어 있고, 마지막으로 $N-1$번 도시와 $1$번 도시가 연결되어 있다. 앞에서 설명한 순서대로 각 도로를 이동하는데 걸리는 시간이 주어질 때, 최단거리가 가장 먼 두 도시 사이의 거리를 출력하라. $N \leq 5 \times 10^5$이고 각 도로를 이동하는데 걸리는 시간은 양의 정수이며, 임의의 두 도시 사이의 최단거리는 $10^9$ 이하이다.
\end{problem}

\begin{problem}
    (10132번) 원형 철길 상의 $N$개의 기차역이 있다. $1 \leq i < N$에 대해 $i$번쨰 지점과 $i+1$번쨰 사이의 거리는 $l_i$, $N$번째 지점과 첫 번째 지점 사이의 거리는 $l_N$이다. 또한, 이 곳에서는 총 $s$대의 기차가 사용된다. 다음은 이 철도 체계에 대한 설명이다.
    \begin{itemize}
        \item $i$번째 기차는 연료를 꽉 치우고 $d_i$만큼 이동할 수 있다. 
        \item 연료는 역에서만 채울 수 있으므로 각 기차는 다음 역에 도달할만큼 충분히 연료가 남아있지 않다면 연료를 채워야 한다. 물론 그 양이 충분하더라도 원하면 연료를 채울 수 있다. 
        \item 각 $s$대의 기차가 하루에 딱 한번 한 기차역에서 출발해 한 번 순환한 후 다시 출발한 기차역으로 돌아온다. 이때 다음 날의 운행을 위해 연료를 꽉 채우고 퇴근해야 한다. 즉, 기차가 운행을 시작할 때에도 연료는 꽉 차 있는 상태이다.
    \end{itemize}
    각 기차에 대해 출발역을 적절히 고르면 도착 이후 마지막 급유를 포함한 기차의 급유 횟수를 최소화할 수 있을 것이다. 각 기차에 대해 가능한 급유 횟수의 최솟값을 출력하라. 만약 어떤 역을 출발역으로 선택하더라도 기차가 역 중간에 연료 부족으로 멈추는 일이 생긴다면 \texttt{NIE}를 출력하라. $N \leq 10^9$, $s \leq 100$, $\sum l_i \leq 10^9$, $1 \leq d_i \leq \sum l_i$이다.
    
    \textit{* 참고: \texttt{NIE}는 폴란드어로 "NO"라는 뜻이다.}
\end{problem}

\begin{problem}
    (14058번) 음 아닌 정수 $K$와 양의 정수 $P$, 그리고 집합 $A = \{a_1, a_2, \cdots, a_N\}$를 생각하자. 모든 $1 \leq i \leq N$에 대해 $0 \leq a_i < P$이고 $a_i$는 모두 다르다. 이제 다음 규칙의 게임을 생각하자.
    \begin{itemize}
        \item 선 플레이어부터 번갈아가며 $A$에서 수 하나를 지운다.
        \item 수를 $K$번 지운 이후 남은 수들의 합이 $P$의 배수이면 선 플레이어가 이기고 아니면 선 플레이어가 진다.
    \end{itemize}
    두 플레이어가 모두 자신이 이길 수 있도록 플레이한다고 가정하였을 때, 총 $T$개의 테스트케이스에 대해 이기는 플레이어를 출력하라. 선 플레이어가 이기면 \texttt{X}를, 후 플레이어가 이기면 \texttt{Y}를 출력하면 된다. $1 \leq K \leq N \leq 5000$, $P \leq 10^{18}$이다.
\end{problem}

\begin{problem} 
    (17030번) 길이 $N$의 수열 $p_i$가 주어진다. 여기서 $p_i$는 확률로 $0 \leq p_i \leq 1$이다. 모든 $l \leq i \leq r$에 대해 각 수를 $p_i$의 확률로 선택할 때 선택되는 수가 \textbf{정확히 하나}인 확률이 최대가 되는 구간 $[l, r]$을 찾아 그때 최대 확률을 출력하라. $N \leq 10^6$이고, 모든 확률 입력 및 출력은 $10^6$배 된 값으로 이루어진다. 즉, 입력되는 수를 $10^6$으로 나눈 것이 실제 $p_i$이고, 출력도 실제 최대 확률에 $10^6$을 곱한 값의 정수부를 출력해야 한다.
\end{problem}

\begin{problem}
    (18888번) $N$개의 서로 다른 점 $P_i$들의 집합 $S$가 주어진다. 임의의 세 점은 한 직선 위에 있지 않다. 네 꼭짓점이 모두 $S$의 원소이며 점 $P_i$기 사각형의 경계를 제외한 내부에 포함되는 사각형의 개수를 $f(P_i)$라고 할 때, $\sum_{i=1}^N {f(P_i)}$를 출력하라. $5 \leq N \leq 2500$이고 $P_i$의 $x$좌표와 $y$좌표는 모두 절댓값이 $10^9$ 이하인 정수이다.
\end{problem}

\begin{problem}
    (20087번**, IOI16) $0$번 역부터 $N-1$번 역까지 잇는 철도가 하나 있다. 모든 $0 \leq i < N-1$에 대해 $i$번 역과 $i+1$번 역 사이를 지나가는데 걸리는 시간은 $l_i$이다. 또한 각 역에는 지선철도로 다른 한 역과 연결되어 있을 수 있다. 만약 연결되어 있다면 $d_i \neq 0$이고 $i$번 지선철도의 길이는 $d_i$이다. $d_i = 0$은 지선철도가 없음을 나타낸다. 여러분은 \textbf{지선철도로만 연결된 역이 아닌 번호가 붙어 있는 역} 중 두 개를 골라 $c$의 시간에 두 역 사이를 이동할 수 있는 급행철도 노선을 신설하려고 한다. \textbf{지선철도로만 연결된 번호가 없는 역을 포함}하여 두 역을 고를 떄 두 역 사이 거리의 최댓값을 최소화시킬 수 있는 곳에 급행철도 노선을 신설할 때, 가능한 두 역 사이의 거리의 최댓값의 최솟값을 출력하라.
    \begin{itemize}
        \item 모든 경우 $1 \leq l_i \leq 10^9$, $0 \leq d_i \leq 10^9$, $1 \leq c \leq 10^9$이다.
        \item (서브테스크 1) $2 \leq N \leq 10$
        \item (서브테스크 2) $2 \leq N \leq 100$
        \item (서브테스크 3) $2 \leq N \leq 250$
        \item (서브테스크 4) $2 \leq N \leq 500$
        \item (서브테스크 5) $2 \leq N \leq 3000$
        \item (서브테스크 6) $2 \leq N \leq 10^5$
        \item (서브테스크 7) $2 \leq N \leq 3 \times 10^5$
        \item (서브테스크 8) $2 \leq N \leq 10^6$
    \end{itemize}
\end{problem}

\begin{problem}
    (24261번) 길이 $N$의 수열 $A_i$와 길이 $M$의 수열 $B_i$가 주어진다. 또한, 모든 $1 \leq i \leq N$에 대해 $A_i \in [1, M]$이고 모든 $1 \leq j \leq M$에 대해 $B_j \in [1, N]$가 성립한다. $A_i$에서 한 개 이상의 원소를 골라 더한 합과 $B_i$에서 한 개 이상의 원소를 골라 더한 합이 같아지는 경우 하나를 출력하라. 선택되는 원소의 개수와 인덱스를 오름차순으로 각각 출력하면 된다. $N, M \leq 10^6$이다.
\end{problem}

\begin{problem}
    (25442번, IOI22) 길이 $N$의 배열 $a_i$와 집합 $A$가 주어질 때, \textit{최대 빈도}와 \textit{최소 빈도}를 다음과 같이 정의하자.
    \begin{itemize}
        \item 정수 $k$에 대해 \textit{$k$의 빈도}는 집합 $A$의 각 원소 $x$에 대해 $a_x=k$인 $x$의 개수이다.
        \item \textit{최대 빈도}는 모든 $1 \leq k \leq 10^9$에 대해 \textit{$k$의 빈도}중 최댓값이다.
        \item \textit{최소 빈도}는 모든 $1 \leq k \leq 10^9$에 대해 \textit{$k$의 빈도}중 0을 제외했을 때 최솟값이다.
    \end{itemize}
    또한 처음 집합 $A = \phi$이고 다음 세 가지 연산을 원하는 순서로 여러 번 사용할 수 있다. 시행 결과는 누적된다.
    \begin{itemize}
        \item $1 \leq i \leq N$인 $i$를 하나 골라 $A$에서 $i$가 있다면 $i$를 제거한다.
        \item $1 \leq i \leq N$인 $i$를 하나 골라 $A$에서 $i$가 없다면 $i$를 추가한다.
        \item $A$의 \textit{최대 빈도}를 알아낸다.
    \end{itemize}
    세 연산을 적절히 사용하여 집합 $\{ 1, 2, \cdots, N \}$의 \textit{최소 빈도}를 알아내어라.
    \begin{itemize}
        \item 모든 경우에서 $1 \leq a_i \leq 10^9$이다.
        \item (서브테스크 1) $N \leq 200$, 세 연산 각각을 $40000$번 이하로 사용
        \item (서브테스크 2) $N \leq 1000$, 세 연산 각각을 $40000$번 이하로 사용
        \item (서브테스크 3) $N \leq 2000$, 세 연산 각각을 $3N$번 이하로 사용
    \end{itemize}
\end{problem}

\begin{problem}
    (26104번) $N$개의 서로 다른 점 $P_i$들의 집합 $S$가 주어진다. 임의의 세 점은 한 직선 위에 있지 않다. 네 꼭짓점이 모두 $S$의 원소이며 그 어떤 점 $P_i$도 사각형의 경계를 제외한 내부에 포함되지 않는 사각형의 개수를 출력하라. $1 \leq N \leq 300$이고 $P_i$의 $x$좌표와 $y$좌표는 모두 절댓값이 $10^9$ 이하인 정수이다.
\end{problem}